\documentclass[11pt]{article}

\usepackage[margin=0.9in]{geometry}
\usepackage[T1]{fontenc}
\usepackage[utf8]{inputenc}
\usepackage{amsmath,amssymb,amsfonts,mathtools,bm}
\usepackage{booktabs,longtable,array}
\usepackage{xcolor}
\usepackage{enumitem}
\usepackage{hyperref}

\hypersetup{colorlinks=true,linkcolor=blue!50!black,urlcolor=blue!50!black}
\setlist{leftmargin=2em}
\setlength{\parskip}{0.35em}
\setlength{\parindent}{0pt}

\newcommand{\C}{\mathbb C}
\newcommand{\R}{\mathbb R}
\newcommand{\Hh}{^{\mathsf H}}
\newcommand{\T}{^{\mathsf T}}
\newcommand{\norm}[1]{\left\lVert #1 \right\rVert}
\newcommand{\abs}[1]{\left\lvert #1 \right\rvert}
\newcommand{\diag}{\operatorname{diag}}
\newcommand{\Ree}{\operatorname{Re}}
\newcommand{\argmin}{\operatorname*{arg\,min}}

\title{Manual Gradients for Code-Constrained OFDM Receiver Ideas\\
\large How the Underwater OFDM Paper Set Fits a Hand-Derived SC-JUNA Update}
\author{Project SONIQUE}
\date{\today}

\begin{document}
\maketitle

\begin{abstract}
This note explains how to use manual gradients, rather than production
automatic differentiation, on the receiver ideas represented by the OFDM
papers in this directory.  The central recommendation is selective: use a
hand-derived gradient for the real coded-bit latent vector, use closed-form
least squares or deterministic candidates for most physical variables, and use
the FEC code constraints as a differentiable prior inside the OFDM receiver.
\end{abstract}

\section{Core Principle}

The suitable common abstraction is
\begin{equation}
  \min_{\bm z,\Theta,\bm u}
  \mathcal L(\bm z,\Theta,\bm u)
  =
  \mathcal D\big(\bm y;\bm s(\bm z),\Theta,\bm u\big)
  +\mathcal R_{\rm phys}(\Theta,\bm u)
  +\lambda_p \mathcal R_{\rm code}(\bm z).
  \label{eq:master}
\end{equation}
Here $\bm z\in\R^{n_b}$ is a real latent vector for coded bits, $\Theta$
contains receiver physics such as Doppler, ICI span, channel coefficients, and
partial-FFT combiners, and $\bm u$ contains nuisance terms such as impulsive
noise or external interference.

The important design choice is:
\begin{quote}
\textbf{Differentiate only the part that benefits from differentiability: the
soft coded-bit latent vector and a few low-dimensional continuous refinements.
Keep support selection, candidate selection, and most channel/interference
updates explicit.}
\end{quote}

\section{Bit Latents and Code Constraint}

For binary coded QPSK, define
\begin{equation}
  \xi_i=\tanh z_i,\qquad \xi_i\in(-1,1),
\end{equation}
and for data carrier $k$ with coded-bit indices $(i_I(k),i_Q(k))$,
\begin{equation}
  s_k(\bm z)=\frac{1}{\sqrt 2}\big(\xi_{i_I(k)}+j\xi_{i_Q(k)}\big).
  \label{eq:qpsk}
\end{equation}
The derivatives are
\begin{equation}
  \frac{\partial s_k}{\partial z_i}
  =
  \begin{cases}
    (1-\xi_i^2)/\sqrt 2, & i=i_I(k),\\
    j(1-\xi_i^2)/\sqrt 2, & i=i_Q(k),\\
    0, & \text{otherwise.}
  \end{cases}
  \label{eq:symbol_derivative}
\end{equation}

For an LDPC parity-check graph with checks $\mathcal N_j$, use the relaxed
parity residual
\begin{equation}
  \mathcal R_{\rm code}(\bm z)
  =
  \sum_j
  \left(1-\prod_{i\in\mathcal N_j}\xi_i\right)^2.
  \label{eq:parity_penalty}
\end{equation}
Let $P_j=\prod_{i\in\mathcal N_j}\xi_i$ and
$P_{j\setminus i}=\prod_{r\in\mathcal N_j\setminus i}\xi_r$.  Then
\begin{equation}
  \frac{\partial \mathcal R_{\rm code}}{\partial z_i}
  =
  -2(1-\xi_i^2)
  \sum_{j:i\in\mathcal N_j}
  (1-P_j)P_{j\setminus i}.
  \label{eq:parity_gradient}
\end{equation}
Do not compute $P_{j\setminus i}$ by division $P_j/\xi_i$; compute the
excluded product directly or with prefix/suffix products so zeros are safe.

\section{Universal Manual Gradient}

Most paper ideas can be written locally as a linear observation model around
the current physical variables:
\begin{equation}
  \widehat{\bm y}_a
  =
  \sum_{k\in\mathcal D_a}\bm a_{ak}(\Theta)s_k(\bm z)
  +\bm B_a(\Theta)\bm u,
  \qquad
  \bm e_a=\widehat{\bm y}_a-\bm y_a .
  \label{eq:local_model}
\end{equation}
For
\begin{equation}
  \mathcal D_{\rm obs}=\sum_a \norm{\bm e_a}_2^2,
\end{equation}
the hand gradient is
\begin{equation}
  \frac{\partial \mathcal D_{\rm obs}}{\partial z_i}
  =
  2\sum_{a:\,k(i)\in\mathcal D_a}
  \Ree\left[
    \bm e_a\Hh \bm a_{a,k(i)}
    \frac{\partial s_{k(i)}}{\partial z_i}
  \right],
  \label{eq:universal_z_gradient}
\end{equation}
where $k(i)$ is the carrier containing bit $i$.  The full $z$ gradient is
\begin{equation}
  \nabla_z\mathcal L
  =
  \nabla_z\mathcal D_{\rm obs}
  +\lambda_p\nabla_z\mathcal R_{\rm code}
  +\nabla_z\mathcal R_{\rm aux}.
  \label{eq:full_z_gradient}
\end{equation}
This is the reusable manual-gradient hook for the whole paper set.

\section{Suitable Concepts From the Papers}

\begin{longtable}{@{}>{\raggedright\arraybackslash}p{0.24\linewidth}>{\raggedright\arraybackslash}p{0.28\linewidth}>{\raggedright\arraybackslash}p{0.38\linewidth}@{}}
\toprule
Paper concept & What to differentiate manually & What should stay explicit \\
\midrule
Partial FFT demodulation
& Differentiate the coded-bit latents through the partial-FFT observation model
$\bm Y_k\approx\sum_q \bm c_{kq}s_q(\bm z)$.  Use
Eq.~\eqref{eq:universal_z_gradient}.
& Partial-FFT windows, the number of partial-FFT views $M$, and the combiner weights should be
deterministic candidates or least-squares updates. \\

Progressive ICI equalization
& Differentiate $\bm z$ for the current ICI half-bandwidth
$K_{\rm ICI}$.  Soft code constraints replace hard tentative decisions.
& Increasing $K_{\rm ICI}$ is a model-selection step.  Use CRC, BP success,
parity residual, and residual energy, not gradient descent on integer span. \\

Nonuniform Doppler compensation
& After a Doppler candidate is chosen, differentiate $\bm z$ through the
resulting OFDM/partial-FFT observation.  Optional v2: hand-derive one scalar
gradient for residual Doppler if interpolation is fixed.
& Coarse resampling, timing, and Doppler search should be a deterministic grid
or preamble-based estimator.  Do not differentiate the full acquisition path in
v1. \\

Sparse delay-Doppler channel estimation
& Differentiate $\bm z$ after a sparse support or dictionary slice has been
chosen.  Reliable soft symbols become code-constrained anchors.
& OMP support selection, basis-pursuit active sets, and dictionary pruning are
not good gradient targets.  Use OMP/BP/proximal steps and then LS. \\

Clustered channel adaptation
& Possible manual gradients for low-dimensional cluster delay/Doppler offsets:
$\partial L/\partial\theta_c=2\Ree[\bm e\Hh(\partial A/\partial\theta_c)\bm s]$.
& Cluster assignment and large support changes should be discrete candidates.
Cluster amplitudes should be LS. \\

Impulsive noise mitigation
& Differentiate $\bm z$ after subtracting the current impulse estimate.  If a
soft impulse mask is used, its scalar gradient can be derived explicitly.
& Detect impulse positions by threshold/GLRT/null-tone residuals.  Estimate
impulse samples by LS or sparse proximal updates. \\

Partial-band partial-block interference
& Differentiate $\bm z$ through
$\widehat{\bm y}=A\bm s(\bm z)+B_{\rm int}\alpha$.  The code penalty prevents
over-cancellation of valid signal energy.
& Interference support in time/frequency should come from GLRT or candidate
sets.  Interference coefficients $\alpha$ should be LS. \\

Nonbinary LDPC coding
& v1: binary-image bit latents and binary parity gradients.  v2: group bit
latents into GF($2^m$) symbols and derive a nonbinary check surrogate.
& Nonbinary BP messages and field operations should remain explicit.  Do not
hide the finite-field algebra inside a black-box differentiator. \\

Asynchronous multiuser OFDM
& Give each user its own latent vector $\bm z^{(u)}$ and parity penalty.  The
gradient for user $u$ uses that user's mixing matrix
$A_u(\tau_u,\Theta_u)$.
& User timing offsets, overlap segmentation, and cancellation order should be
candidates or small-grid searches. \\

Physical-layer network coding
& Use manual gradients for separate latents $(\bm z_A,\bm z_B)$, network-coded
latents $\bm z_\oplus$, or composite latents with an XOR consistency penalty.
& Which PLNC decoding mode to trust is a candidate-scoring problem, not a
continuous-gradient problem. \\

Adaptive modulation/coding and outer erasure coding
& No packet-level gradient is needed.  Use manual-gradient receiver outputs as
features: parity residual, LLR margin, selected $M$, selected $K_{\rm ICI}$,
CRC status, and residual energy.
& Mode selection and erasure-code allocation are system-layer policies. \\
\bottomrule
\end{longtable}

\section{Concrete Losses for the Most Useful Modules}

\subsection{Partial FFT and ICI}

Let $\bm Y_k\in\C^M$ be the partial-FFT vector for carrier $k$, and let
$\bm c_{kq}\in\C^M$ be the response from transmitted carrier $q$ across
the $M$ partial-FFT views at received carrier $k$.  With
$\mathcal B_K(k)=\{q:\abs{q-k}\le K\}$,
\begin{equation}
  \bm r_k
  =
  \sum_{q\in\mathcal B_K(k)}\bm c_{kq}s_q(\bm z)-\bm Y_k.
\end{equation}
Use
\begin{equation}
  \mathcal L_{\rm pfft}
  =
  \lambda_d\sum_k\norm{\bm r_k}_2^2
  +\lambda_p\mathcal R_{\rm code}(\bm z).
\end{equation}
Then for bit $i$ on carrier $q$,
\begin{equation}
  \frac{\partial \mathcal L_{\rm pfft}}{\partial z_i}
  =
  2\lambda_d
  \sum_{k:q\in\mathcal B_K(k)}
  \Ree\left[
    \bm r_k\Hh\bm c_{kq}\frac{\partial s_q}{\partial z_i}
  \right]
  +\lambda_p
  \frac{\partial\mathcal R_{\rm code}}{\partial z_i}.
  \label{eq:pfft_gradient}
\end{equation}

\subsection{Combiner Symbol Tie}

If $\widehat s_k=\bm w_k\Hh\bm Y_k$, add
\begin{equation}
  \mathcal L_{\rm tie}
  =
  \rho\sum_{k\in\mathcal D}\abs{\widehat s_k-s_k(\bm z)}^2 .
\end{equation}
For bit $i$ on carrier $k$,
\begin{equation}
  \frac{\partial \mathcal L_{\rm tie}}{\partial z_i}
  =
  -2\rho\Ree\left[
    \big(\widehat s_k-s_k\big)^*
    \frac{\partial s_k}{\partial z_i}
  \right].
  \label{eq:tie_gradient}
\end{equation}
The negative sign appears because the residual is $\widehat s_k-s_k$.

\subsection{Sparse or Clustered Channel Refinement}

After support selection, write the model as
\begin{equation}
  \widehat{\bm y}=A(\bm\theta)\bm s(\bm z),\qquad
  \bm e=\widehat{\bm y}-\bm y.
\end{equation}
For a scalar cluster or path parameter $\theta_c$,
\begin{equation}
  \frac{\partial}{\partial\theta_c}\norm{\bm e}^2
  =
  2\Ree\left[
    \bm e\Hh
    \frac{\partial A(\bm\theta)}{\partial\theta_c}
    \bm s(\bm z)
  \right].
  \label{eq:theta_gradient}
\end{equation}
Use this only for a small number of continuous refinements, such as residual
Doppler offset or cluster delay drift.  If $\partial A/\partial\theta_c$ is
messy or discontinuous because of interpolation boundaries, prefer a local grid
search.

\subsection{Impulse and External-Interference Nuisance Terms}

For impulsive noise or partial-band interference,
\begin{equation}
  \widehat{\bm y}=A\bm s(\bm z)+B\bm u,\qquad
  \bm e=\widehat{\bm y}-\bm y.
\end{equation}
Given $\bm z$, update $\bm u$ by
\begin{equation}
  \bm u
  =
  \argmin_{\bm u}\norm{A\bm s(\bm z)+B\bm u-\bm y}^2
  +\lambda_u\mathcal R_u(\bm u),
\end{equation}
using LS if $\mathcal R_u=0$, or a sparse proximal step if $\mathcal R_u$ is an
$\ell_1$ penalty.  Given $\bm u$, update $\bm z$ by
Eq.~\eqref{eq:universal_z_gradient}.  This alternating structure is safer than
trying to manually differentiate every nuisance detector.

\section{Recommended Solver Schedule}

\begin{enumerate}
  \item Generate deterministic candidates for coarse Doppler, partial-FFT
  number of partial-FFT views $M$, and ICI span $K_{\rm ICI}$.
  \item For each candidate, initialize channel and response coefficients from pilots
  and optional sparse support selection.
  \item Alternate:
  \begin{enumerate}
    \item update channel and interference variables by LS, grid, or proximal
    rule;
    \item update $\bm z$ by a few manual-gradient steps using
    Eq.~\eqref{eq:full_z_gradient};
    \item run BP or use BP extrinsic information to refresh LLRs and confidence.
  \end{enumerate}
  \item Score candidates by pilot residual, data residual, parity residual,
  LLR margin, BP success, and CRC.
  \item Treat high parity residual or CRC failure as an erasure signal for any
  outer erasure code or AMC policy.
\end{enumerate}

\section{What Not to Differentiate}

Manual gradients are most useful when the mapping is smooth and low-dimensional.
They are not the right tool for every idea in the papers.

\begin{itemize}
  \item Do not gradient-descent integer quantities: ICI span, number of partial-FFT views,
  pilot pattern, LDPC graph, support size, or modulation mode.
  \item Do not differentiate through hard OMP support selection, GLRT decisions,
  CRC decisions, or BP stopping rules.
  \item Do not make the finite-field LDPC algebra an opaque real-valued neural
  layer.  Keep BP explicit and use the relaxed parity penalty only as a
  receiver-side soft constraint.
  \item Do not let impulse/interference nuisance variables explain everything.
  Always score with parity and CRC so the receiver prefers code-valid signal
  explanations.
\end{itemize}

\section{Summary}

The best manual-gradient use of the paper set is not a fully differentiable
receiver.  It is a structured alternating receiver:
\begin{equation}
  \boxed{
  \begin{gathered}
  \text{LS/grid/prox for physics}\\
  {}+\text{ manual gradient for coded bits}\\
  {}+\text{ BP/CRC/parity for code consistency}
  \end{gathered}
  }
\end{equation}
Partial FFT, progressive ICI, Doppler compensation, sparse and clustered
channels, interference cancellation, multiuser reception, and PLNC can all plug
into this pattern by changing the observation matrix or nuisance basis while
leaving the hand-derived coded-bit gradient intact.

\end{document}
